% 【本文件写什么】impl 实现层：virtio-mmio
% - 定位：VirtIO MMIO 传输，RISC-V 主线块/网。
% - crate 路径：os/components/wateros-driver/driver-block/block-impl/impl-virtio-mmio/src/
% - 如何实现对应 api-v0 trait：关键类型、状态机、数据结构
% - 算法或硬件交互流程，可用伪代码/流程图
% - 本 impl 启用的 Cargo [features] 与 arch feature，impl-riscv64 等
% - 与其它 impl 变体的差异、切换 feature 时的影响
% - 性能/内存假设与已知限制
% 【事实来源】
% - 上述 crate 源码与 Cargo.toml
% - os/feature-tree.txt 中指向本 impl 的 feature 链

\subsubsection{impl-virtio-mmio}

\paragraph{启用链。}由根\code{wateros-driver}的\code{impl-qemu-riscv64-opensbi}打开\code{block/impl-virtio-mmio}。平台impl在DTB节点\code{compatible=="virtio,mmio"}且\code{device\_id==2}时调用\code{VirtioBlkDevice::from\_mmio}。

\paragraph{HAL与DMA。}\code{virtio-drivers}队列缓冲经\code{VirtioMmioHal::dma\_alloc}向\code{wateros-mm-frame-alloctor}要物理连续、页对齐、已清零帧；恒等映射下\code{paddr==vaddr}。\code{dma\_dealloc}按页号归还帧池。须在\code{init\_frame\_allocator}之后实例化设备，否则分配失败映射为\code{DriverError::Unsupported}。

\paragraph{读写路径。}\code{read\_blocks}/\code{write\_blocks}将\code{Lba}转为\code{usize}后委托\code{VirtIOBlk}；传输错误折叠为\code{DriverError::IoError}。大块写，$\geq$4096字节会打\code{trace}日志便于iozone诊断。

\begin{rustcode}
pub struct VirtioBlkDevice {
    inner: VirtIOBlk<VirtioMmioHal, MmioTransport<'static>>,
}

impl VirtioBlkDevice {
    pub fn from_mmio(mmio: MmioRegion) -> DriverResult<Self> {
        let header = NonNull::new(mmio.base as *mut VirtIOHeader)
            .ok_or(DriverError::InvalidDtb)?;
        let transport = unsafe { MmioTransport::new(header, mmio.size) }
            .map_err(|_| DriverError::Unsupported)?;
        let inner = VirtIOBlk::<VirtioMmioHal, MmioTransport>::new(transport)
            .map_err(|_| DriverError::Unsupported)?;
        Ok(Self { inner })
    }
}
\end{rustcode}

\paragraph{限制。}仅使用DTB\code{reg}首段MMIO窗口；未挂接\code{DeviceInfo.irq}对应的PLIC处理，依赖轮询或virtio内部完成传输。与PCI路径相比无BAR分配步骤，适合QEMU\code{virt}上\code{virtio-blk-device}节点。
